\documentclass[11pt]{article}
\usepackage{fontspec}
\usepackage[french,english]{babel}
\usepackage{amsmath}
\usepackage{amssymb}
\usepackage{amsthm}
\usepackage{mathrsfs}
\usepackage{tikz-cd}
\usepackage{enumitem}

\newcommand{\printedpage}[1]{\marginpar{\footnotesize p.~#1}}
\newcommand{\msp}[1]{\marginpar{\footnotesize [#1]}}
\newcommand{\uHom}{\underline{\mathrm{Hom}}}
\newcommand{\uGrp}{\underline{\mathrm{Grp}}}
\newcommand{\uuGrp}{\underline{\underline{\mathrm{Grp}}}}
\newcommand{\fl}{\longrightarrow}

\begin{document}

\printedpage{4}
% head: A. Grothendieck : Dérivateurs, Ch. VIII, Édition des Universités de Montpellier II et Paris VII
% foot: Version 2 février 2002

\textbf{Theorem.} \textit{Let $f : S \fl T$ be an arrow in} (Cat)\textit{, $G$ a stack in groupoids over $S$, $H$ a stack in groupoids over $T$. We then have a canonical equivalence of categories}
\[
\boxed{\ \uHom_{\uGrp(S)}(G, f^*H) \ \simeq\ \uHom_{\uGrp(T)}(f_!G, H).\ }
\]

The adjunction arrows
\[
\left\{
\begin{array}{lclcrcl}
f_!f^*H & \fl & H & \qquad\qquad & H & \in & \mathrm{Ob}\,\uGrp(T)\\[4pt]
f^*f_!G & \longleftarrow & G & \qquad\qquad & G & \in & \mathrm{Ob}\,\uGrp(S)
\end{array}
\right.
\]
should be definable without any problem. For the second, we have to define for $s \in S$
\[
G_s \ \fl\ (f^*(f_!(G)))_s = (f_!(G))_{f(s)} \overset{\text{def}}{=} \Pi(s\backslash G),
\]
now we have
\[
G_s \fl s\backslash G,
\]
hence
\[
\underbrace{\Pi(G_s)}_{\begin{subarray}{l} =\ G_s\\ \text{since } G_s\\ \text{is a}\\ \text{groupoid}\end{subarray}} \ \fl \Pi(s\backslash G),
\]
O.K. This is functorial in $s$ \dots

\msp{page 5}
For the arrow $f_!f^*H \fl H$ over $T$, we have to find for every $t \in \mathrm{Ob}\,T$, and functorially in $t$,
\[
\underbrace{(f_!f^*H)_t}_{\overset{\text{def}}{=}\Pi(t\backslash f^*(H))} \quad \fl \quad \underbrace{H_t}_{=\,\Pi(H_t)} \ ,
\]
it would therefore suffice to define
\[
t\backslash f^*(H) \ \fl\ H_t,
\]
but we have the cartesian diagram
\[
\begin{tikzcd}[column sep=large, row sep=large]
f^*H = H \times_S T \arrow[r] \arrow[d] & H \arrow[d] \\
S \arrow[r, "f"] & T \ni t
\end{tikzcd}
\]
\texttt{[rather} $H \times_T S$\texttt{]} inducing
\[
f^*H \ \fl\ H,
\]
\printedpage{5}
% head: A. Grothendieck : Dérivateurs, Ch. VIII, Édition des Universités de Montpellier II et Paris VII
% foot: Version 2 février 2002
hence
\[
t\backslash f^*H \ \fl\ t\backslash H \qquad\quad \text{(instead of mapping into $H_t$),}
\]
but on the other hand, since $H$ is fibered over $T$, hence smooth over $T$, the inclusion
\[
H_t \fl t\backslash H
\]
is a Hot-equivalence. Or rather, let us say: $f^*H \fl H$ is a morphism of categories over $T$, inducing
\[
\underbrace{\Pi(f^*H/T)}_{\overset{\text{def}}{=}\ f_!(f^*H)} \ \fl\ \Pi(H/T),
\]

\msp{page 6}
now the canonical $T$-functor
\[
H \fl \Pi(H/T) \qquad\quad (= \text{category over } T \text{ defined by } t \longmapsto \Pi(t\backslash H))
\]
is an equivalence, because $H$ is a groupoid fibration over $T$. So we can find a $T$-quasi-inverse
\[
\Pi(H/T) \fl H
\]
and by composing we find
\[
f_!(f^*H) \ \fl\ \Pi(H/T) \ \fl\ H
\]
the arrow sought.

Having the adjunction 2-arrows, I stop there with the verifications!

\printedpage{6}
% head: A. Grothendieck : Dérivateurs, Ch. VIII, Édition des Universités de Montpellier II et Paris VII
% foot: Version 2 février 2002

\section*{2\quad Base change questions}

\[
\begin{tikzcd}[column sep=large, row sep=large]
S & S' \arrow[l, "q"'] \arrow[d, "f'"] \\
T \arrow[from=1-1, "f"'] & T' \arrow[l, "p"]
\arrow[from=2-2, to=1-1, bend right=50, shorten <=18pt, shorten >=18pt, "\gamma"']
\end{tikzcd}
\qquad\qquad\qquad
p^*f_!(G) \longleftarrow f'_!q^*(G)
\tag*{$(*)$}
\]

The arrow is defined, provided we have commutativity or a commutativity arrow $\gamma : pf' \fl fq$ (not necessarily \texttt{[an]} isomorphism). By adjunction, it amounts to the same thing to give
\[
\begin{tikzcd}[row sep=large]
p_!f'_!q^*G \arrow[r] \arrow[d, dash, "\wr"', "\text{transitivity}"] & f_!(G) \\
(pf')_!q^*G \arrow[d, "\gamma"] & \\
(fq)_!q^*G \arrow[d, dash, "\text{transitivity}\ \wr"'] & \\
f_!(q_!q^*)G. \arrow[uuur, dashed, "\text{since } q_!q^*G \longrightarrow G"'] &
\end{tikzcd}
\]

We wonder when the arrow is \texttt{[an]} isomorphism.

\msp{page 7}
The following theorem must be true, and follows from my general \textit{fourbis} (gear).

% vertical bar printed in the left margin beside the following theorem
\textbf{Theorem.} \textit{If the square $(*)$ is} cartesian \textit{in} (Cat)\textit{, $\gamma$ an isomorphism, and if we assume that $f$ is} 1\textit{-smooth, or $p$} 1\textit{-proper, then the $T'$-functor of base change is an equivalence of $T'$-categories.}

The general ``\textit{fourbis}'' would be in terms of an ad hoc derivator $\mathbf{D}$, that of relative homotopy 1-types. I will return to this later.

We ask the question when the 2-functor associated to $f : S \fl T$
\[
f^* :\ \uuGrp(T) \ \fl\ \uuGrp(S)
\]
\printedpage{7}
% head: A. Grothendieck : Dérivateurs, Ch. VIII, Édition des Universités de Montpellier II et Paris VII
% foot: Version 2 février 2002
is 0-faithful, 1-faithful, 2-faithful (\footnote[2]{The $i$-faithfulness of $f^*$ ($i \in \{0,1,2\}$) is equivalent to that of the functors (obtained by taking $G = e_T$)
\[
\underline{\Gamma}(H/T) \ \fl\ \underline{\Gamma}(f^*H/S).
\]}). If $G$, $H$ are groupoid fibrations over $T$, we consider
\[
\uHom_{\uuGrp(T)}(G,H) \ \fl\ \uHom_{\uuGrp(S)}(f^*G, f^*H),
\]
and according as this functor is always faithful, or fully faithful, or an equivalence of categories, we say that $f^*$ is 0-, 1-, or 2-faithful. In the 2-faithful case, this implies that if $K$ is a groupoid fibration over $S$, if it is equivalent to an $f^*G$, then $G$ is determined up to $T$-equivalence, itself determined up to unique isomorphism.

\msp{page 8}
Let $i \in \mathbf{N}$. We say that $f$ is $i$-aspheric if $\forall t \in T$, $Z = S/t$ is $i$-aspheric, i.e.
\[
\pi_0(Z),\ \dots,\ \pi_i(Z) \qquad\quad \text{trivial.}
\tag*{$(*)$}
\]

We also define $(-1)$-aspheric: the $S/t$ are nonempty. Likewise for the notion of $i$-coaspheric.

In fact, for $i \in [-1,+\infty[ \ \subseteq \mathbf{Z}$, we define a derivator $\mathbf{D}_i$ of relative homotopy $i$-types (cf. below). Hence the notions of a $\mathbf{D}_i$-aspheric object ($=$ $i$-aspheric), of a $\mathbf{D}_i$-equivalence, of a $\mathbf{D}_i$-aspheric and $\mathbf{D}_i$-coaspheric functor, of a $\mathbf{D}_i$-smooth and $\mathbf{D}_i$-proper functor, of a $\mathbf{D}_i$-fibration and of a $\mathbf{D}_i$-fibration that is \textit{trivial} (i.e. with $\mathbf{D}_i$-aspheric fibers). We should have the

\textbf{Theorem 3.} \textit{Let $f : S \fl T$ be in} (Cat)\textit{, $i \in \{0,1,2\}$. Suppose one of the following conditions is satisfied:}

\begin{itemize}[leftmargin=3em]
\item[\textit{a)}] \textit{$f$ is $(i-1)$-proper, and $(i-1)$-aspheric (is expressed by the condition that certain categories, such as the $S/t$, be $(i-1)$-aspheric).}
\item[\textit{b)}] \textit{$f$ is $(i-1)$-smooth, and $(i-1)$-coaspheric.}
\end{itemize}

\textit{Then:}

\begin{itemize}[leftmargin=3em]
\item[\textit{1\textsuperscript{o})}] \textit{The inverse image functor $f^*$}
\[
f^* : \uGrp(T) \ \fl\ \uGrp(S)
\]
\textit{is $i$-faithful.}
\item[\textit{2\textsuperscript{o})}] \textit{In the case $i = 2$ ($f$ fully faithful), we}
\printedpage{8}
% head: A. Grothendieck : Dérivateurs, Ch. VIII, Édition des Universités de Montpellier II et Paris VII
% foot: Version 2 février 2002
\msp{page 9}
\textit{also have this: if $G \in \mathrm{Ob}\,\uGrp(S)$ is a groupoid fibration over $S$, then $G$ is equivalent to a groupoid fibration of the form $f^*(H)$ if and only if for every $t \in T$, the restriction $G_t$ of $G$ to $S_t$ is equivalent to a constant groupoid fibration over $S_t$. (}\textbf{NB} \textit{Condition a) or b) already implies that $S_t$ is} 1\textit{-connected, i.e. that its fundamental groupoid is equivalent to $e$.)}
\item[\textit{3\textsuperscript{o})}] \textit{Let $Z$ be an object} 1\textit{-connected in} (Cat)\textit{, $G$ a groupoid fibration over $Z$. For $G$ to be $Z$-equivalent to a constant groupoid over $Z$, it is necessary and sufficient that for every arrow $u$ in $Z$, the morphism of groupoids $u^*$ induced on the fibers of $G$ be a $Z$-equivalence} (\footnote[3]{We then say that $G$ is \textit{essentially locally constant} over $Z$.}).
\end{itemize}

\textbf{Corollary.} \textit{The necessary and sufficient condition of 2\textsuperscript{o}) can also be written as follows: for every arrow $u$ of $S$ ``vertical'' over $T$, i.e. such that $f(u)$ is an identity arrow, the functor $u^*$ induced by $u$ on the fibers of $G$ be an equivalence of groupoids.}

\end{document}
